
\section{哈密顿\trDiffeomorphism{}群 \texorpdfstring{$\mathcal{D}$}{D} 上的%
特殊\trMetric{} \texorpdfstring{$d$}{d}}
\label{sec:special-metric-d}

我们考察具有\trStandard\trSymplecticStructure{} $\omega=\omega_{0}$ 的%
\trSymplectic{}空间 $(\R^{2n},\omega)$，并以 $\mathcal{H}$ 表示全体\trSmooth{}%
且\trCompactSupport{}的\trHamiltonianFunction{} $H=H(t,x):[0,1]\times\R^{2n}%
\to\R$ 所成的向量空间。对每个 $H\in\mathcal{H}$，都有 $\R^{2n}$ 上的一个%
\trTimeDependent\trHamiltonianVectorField{} $X_{H}$ 与之相联系，其定义为
\[
i_{X_{H_{t}}}\omega=-dH_{t},\qquad\text{其中}\ H_{t}(x)=H(t,x).
\]

由于\trVectorField{} $X_{H}$ 具有\trCompactSupport，诸\trHamiltonianEquation
\[
\dot{x}=X_{H}(t,x),\qquad x(0)=x_{0}\in\R^{2n}
\]
对每个给定的初值 $x_{0}\in\R^{2n}$ 都可在整个时间区间 $[0,1]$ 上求解。于是，%
对 $t\in[0,1]$ 我们便得到\trSymplectic\trMapping{}的一个单参数族%
$\varphi_{H}^{t}$，它由
\[
\varphi_{H}^{t}(x_{0})=x(t),
\]
定义，其中 $x(t)$ 表示该方程对初值 $x_{0}$ 的解。显然，当 $\lvert x\rvert$%
充分大时 $\varphi_{H}^{t}(x)=x$。我们以
\[
\varphi_{H}=\varphi_{H}^{1}
\]
记\trFlow{} $\varphi_{H}^{t}$ 的\trTimeOneMap，并如此定义紧支撑%
\trHamiltonianDiffeomorphism{}群 $\mathcal{D}$：
\[
\mathcal{D}=\bigl\{\,\varphi_{H}\ \bigm|\ H\in\mathcal{H}\,\bigr\}.
\]

回顾：一个映射 $\varphi$ 的\trSupport{}定义为\trClosed{}集%
$\operatorname{supp}(\varphi)=\operatorname{clos}\bigl\{\,x\ \bigm|\ %
\varphi(x)\neq x\,\bigr\}$。每个紧支撑\trSymplecticDiffeomorphism{}都能按此%
方式由一个\trTimeDependent\trHamiltonianVectorField{}的\trFlow{}插值而得；%
然而，当 $n>2$ 时，该\trHamiltonianFunction{}能否选自 $\mathcal{H}$，尚属%
未知。由 M.~\trGromov{} [107] 的结果可推知：在 $\R^{4}$ 这一特殊情形，群%
$\mathcal{D}$ 是\trContractible{}的。特别地，$\R^{4}$ 中每个紧支撑%
\trSymplecticDiffeomorphism{} $\varphi$ 都可写成 $\varphi=\varphi_{H}$（%
$H\in\mathcal{H}$），从而属于 $\mathcal{D}$。

本章的目标是研究 $\R^{2n}$ 上这个由\trHamiltonian\trMapping{}组成的群%
$\mathcal{D}$ 的若干几何性质。我们先给出两个预备性的形式表述，它们在后面%
将会用到。

\begin{definition}\label{def:function-operations}
给定 $H,K\in\mathcal{H}$ 与 $\vartheta\in\mathcal{D}$，我们如下定义%
$\mathcal{H}$ 中的函数 $\overline{H}$、$H\#K$ 与 $H_{\vartheta}$：
\begin{align*}
\overline{H}(t,x)\ &=\ -H\bigl(t,\varphi_{H}^{t}(x)\bigr)\\
(H\#K)(t,x)\ &=\ H(t,x)+K\bigl(t,\bigl(\varphi_{H}^{t}\bigr)^{-1}(x)\bigr)\\
H_{\vartheta}(t,x)\ &=\ H\bigl(t,\vartheta^{-1}(x)\bigr).
\end{align*}
\end{definition}

\begin{proposition}\label{prop:calculus-formulae}
对 $H,K\in\mathcal{H}$，下列公式成立：
\begin{align*}
\varphi_{\overline{H}}^{t}\ &=\ \bigl(\varphi_{H}^{t}\bigr)^{-1}\\
\varphi_{H\#K}^{t}\ &=\ \varphi_{H}^{t}\circ\varphi_{K}^{t}\\
\vartheta\circ\varphi_{H}^{t}\circ\vartheta^{-1}\ &=\ \varphi_{H_{\vartheta}}^{t}\\
\bigl(\varphi_{H}^{t}\bigr)^{-1}\circ\varphi_{K}^{t}\ &=\ \varphi_{G}^{t}
\end{align*}
其中 $G=\overline{H}\#K=(K-H)(t,\varphi_{H}^{t})$\footnote{原书印作「$(K-H)(t,\varphi_{H}^{t})$」，疑漏自变量 $x$（按前文约定应为 $(K-H)(t,\varphi_{H}^{t}(x))$），此处照录原书面貌。——译者注}。
\end{proposition}

\begin{proof}
第三个公式是\trHamiltonianVectorField{}在\trSymplectic{}变换下的%
\trTransformationLaw，正如引言中所述。第一个公式由第二个公式推得，而最后%
一个公式是前两个公式的推论。为证第二个公式，我们简缩记号并注意到
\begin{align*}
\frac{d}{dt}\varphi^{t}\ &=\ X_{H}\circ\varphi^{t}\quad\text{以及}\quad%
\varphi^{0}=\id\\
\frac{d}{dt}\psi^{t}\ &=\ X_{K}\circ\psi^{t}\quad\text{以及}\quad\psi^{0}=\id.
\end{align*}
我们必须证明 $\varphi^{t}\circ\psi^{t}$ 是 $X_{H\#K}$ 的\trFlow。求导便得
\begin{align*}
\frac{d}{dt}\bigl(\varphi^{t}\circ\psi^{t}\bigr)\ &=\ \Bigl(\frac{d}{dt}%
\varphi^{t}\Bigr)\circ\psi^{t}+\bigl(d\varphi^{t}\circ\psi^{t}\bigr)\cdot%
\frac{d}{dt}\psi^{t}\\
&=\ X_{H}\bigl(\varphi^{t}\circ\psi^{t}\bigr)+\Bigl[d\varphi^{t}\circ%
\bigl(\varphi^{t}\bigr)^{-1}\circ\varphi^{t}\circ\psi^{t}\Bigr]\cdot X_{K}%
\circ\Bigl[\bigl(\varphi^{t}\bigr)^{-1}\circ\varphi^{t}\circ\psi^{t}\Bigr].
\end{align*}
由\trHamiltonianVectorField{}的\trTransformationLaw，第二项等于
\[
X_{K\circ(\varphi^{t})^{-1}}\circ\bigl(\varphi^{t}\circ\psi^{t}\bigr),
\]
论断得证。
\end{proof}

按定义，$\mathcal{D}$ 中的一条\trSmooth{}弧是指一个映射 $t\mapsto\psi^{t}$：%
它把 $I=[0,1]$ 映入 $\mathcal{D}$，使得 $(t,x)\mapsto\psi^{t}(x)$ 是从%
$I\times\R^{2n}$ 到 $\R^{2n}$ 的一个\trSmooth\trMapping，并且存在 $\R^{2n}$%
中的一个\trCompact{}集，它含有诸 $\psi^{t}$（$t\in I$）的全部\trSupport。%
我们将以 $\mathcal{A}$ 表示满足 $\psi^{0}=\id$ 的\trSmooth{}弧%
$[0,1]\to\mathcal{D}:t\mapsto\psi^{t}$ 之集。

\begin{proposition}\label{prop:arc-bijection}
映射 $H\mapsto\psi_{H}^{t}$ 是从 $\mathcal{H}$ 到 $\mathcal{A}$ 上的一个%
\textbf{\trBijection}（bijection）。
\end{proposition}

\begin{proof}
若 $H$ 与 $K$ 对 $t\in I$ 满足 $\psi_{H}^{t}=\psi_{K}^{t}$，则由求导得%
$X_{H}=X_{K}$，从而 $dH_{t}=dK_{t}$；又因这些函数具有\trCompactSupport，故%
$H(t,x)=K(t,x)$。这证明了该映射的单射性；为证其满射性，我们取 $\mathcal{D}$%
中一条满足 $\psi^{0}=\id$ 的\trSmooth{}弧 $t\mapsto\psi^{t}$，并定义%
\trTimeDependent\trVectorField
\[
X_{t}:=\frac{d}{dt}\psi^{t}\circ\bigl(\psi^{t}\bigr)^{-1},
\]
使得
\[
\frac{d}{dt}\psi^{t}=X_{t}\bigl(\psi^{t}\bigr)\quad\text{且}\quad\psi^{0}=\id.
\]
由于 $(\psi^{t})^{*}\omega=\omega$，我们有 $L_{X_{t}}\omega=0$。鉴于%
$d\omega=0$，我们得 $d(i_{X_{t}}\omega)=0$。\trOneForm{} $i_{X_{t}}\omega$%
既然\trClosed{}，由\trPoincare{}引理它在 $\R^{2n}$ 上便\trExact。于是存在%
函数 $H:I\times\R^{2n}\to\R$，由
\[
H(t,x)=-\int_{0}^{1}\omega\bigl(X_{t}(sx),x\bigr)\,ds,
\]
确定，它满足 $-dH_{t}=i_{X_{t}}\omega$。由于 $X_{t}$ 具有\trCompactSupport，%
故对 $t\in I$ 与 $\lvert x\rvert\geq R$（$R$ 为某个大数）有 $dH_{t}(x)=0$，%
并且存在 $K\in\mathcal{H}$ 使 $dH_{t}=dK_{t}$；它由 $H(t,x)-H(t,x^{*})$ 给出，%
其中 $x^{*}$ 满足 $\lvert x^{*}\rvert\geq R$。从而 $X_{t}=X_{K_{t}}$，即所%
要证的。
\end{proof}

接下来，我们对属于 $\mathcal{A}$ 的一条弧 $t\mapsto\psi^{t}$ 定义其长度。%
鉴于\cref{prop:arc-bijection}，存在唯一的 $H\in\mathcal{H}$ 生成这条弧：
\[
\psi^{t}=\varphi_{H}^{t}.
\]
对 $H\in C^{\infty}_{c}\bigl(\R^{2n},\R\bigr)$ 定义\trNorm{}（称为%
\trOscillation）
\begin{equation}\label{eq:osc-norm}
\lVert H\rVert=\sup_{x}H(x)-\inf_{x}H(x).
\end{equation}
则
\[
\lVert H\circ\varphi\rVert=\lVert H\rVert
\]
对 $\R^{2n}$ 的每个\trDiffeomorphism{} $\varphi$ 成立。若现在 $H\in\mathcal{H}$，%
我们便可使\trNorm{} $\lVert X_{H_{t}}\rVert:=\lVert H_{t}\rVert$ 与%
\trHamiltonianVectorField{} $X_{H_{t}}$ 相联系，并发现
\[
\int_{0}^{1}\bigl\lVert\dot{\psi}^{t}\bigr\rVert\,dt=\int_{0}^{1}%
\bigl\lVert X_{H_{t}}\bigl(\psi^{t}\bigr)\bigr\rVert\,dt=\int_{0}^{1}%
\lVert H_{t}\rVert\,dt.
\]
通过在 $\mathcal{H}$ 中引入\trNorm
\begin{equation}\label{eq:H-norm-int}
\lVert H\rVert=\int_{0}^{1}\lVert H_{t}\rVert\,dt=\int_{0}^{1}%
\Bigl[\sup_{x}H(t,x)-\inf_{x}H(t,x)\Bigr]\,dt,
\end{equation}
它在 $H$ 为自治时与定义 \eqref{eq:osc-norm} 相合，我们便把连接 $\psi^{1}$ 与%
$\id$ 的弧 $\psi^{t}=\varphi_{H}^{t}$ 的长度定义为
\begin{equation}\label{eq:arc-length}
l\bigl(\psi^{[0,1]}\bigr)=\lVert H\rVert.
\end{equation}

\begin{definition}\label{def:energy}
$\mathcal{D}$ 中一个\trSymplectic\trMapping{} $\psi$ 的\trEnergy{}是数%
$E(\psi)\geq0$，定义为
\[
E(\psi)=\inf\bigl\{\,l\bigl(\psi^{[0,1]}\bigr)\ \Bigm|\ \mathcal{A}\text{ 中的弧}\ %
\psi^{t}\ \text{满足}\ \psi^{1}=\psi\,\bigr\}.
\]
\end{definition}

显然
\begin{equation}\label{eq:energy-le-norm}
E(\varphi_{H})\leq\lVert H\rVert.
\end{equation}
$\mathcal{H}$ 中的\trHamiltonianFunction{} $H=0$ 生成\trIdentityMap{} $\psi=\id$，%
从而 $E(\id)=0$。反之，
\begin{equation}\label{eq:energy-positive}
E(\psi)>0,\quad\text{若}\ \psi\neq\id.
\end{equation}
但这一关键的事实绝非显然。它将由下述关于\trEnergy{}函数的\trEstimate{}推出，其中%
$c_{0}(U)=c_{0}(U,\omega)$ 表示第 3 章中引入的那个特殊\trCapacity。为记号%
方便计，我们令 $c_{0}(\emptyset)=0$，并回顾 $Z(R)$ 表示半径为 $R$ 的%
\trSymplectic\trCylinder。

\begin{theorem}[\trEnergy-\trCapacity{}不等式]\label{thm:energy-capacity-inequality}
对每个 $\psi\in\mathcal{D}$，
\[
\sup\bigl\{\,c_{0}(U)\ \Bigm|\ U\ \text{开且}\ \psi(U)\cap U=\varnothing\,\bigr\}%
\ \leq\ E(\psi)
\]
\[
E(\psi)\ \leq\ 16\inf\bigl\{\,\pi R^{2}\ \Bigm|\ \text{存在}\ \varphi\in\mathcal{D}%
\ \text{使}\ \operatorname{supp}\bigl(\varphi\psi\varphi^{-1}\bigr)\subset Z(R)\,\bigr\}.
\]
\end{theorem}

特别地，我们看到：\trEnergy{} $E(\psi)$ 总大于每个能被映射 $\psi$ 从其自身%
挪开的\trOpenSet{} $U$ 的\trCapacity{} $c_{0}(U)$。现在容易推出论断%
\eqref{eq:energy-positive}：的确，取一个满足 $\psi(x^{*})\neq x^{*}$ 的点%
$x^{*}$，便存在半径为 $\varepsilon>0$ 的开\trBall{} $B(x^{*},\varepsilon)=U$，%
它从其自身被挪开，即 $\psi(U)\cap U=\varnothing$。利用 $c_{0}(U)=\pi\varepsilon^{2}$，%
我们便从\cref{thm:energy-capacity-inequality}得出%
$E(\psi)\geq\pi\varepsilon^{2}>0$，论断 \eqref{eq:energy-positive} 得证。

\cref{thm:energy-capacity-inequality} 中第一个\trEstimate{}的证明将基于%
\trActionFunctional{}的一个特定\trCriticalPoint{}的存在性，它将在下文的 5.5%
与 5.6 两节中进行。

\begin{theorem}[$C^{0}$-\trEnergy{}估计]\label{thm:c0-energy-estimate}
对每个 $\psi\in\mathcal{D}$，
\[
E(\psi)\ \leq\ C\ \operatorname{diameter\ supp}\,(\psi)\ %
\bigl\lvert\id-\psi\bigr\rvert_{C^{0}}\footnote{原书下标印作 $C^{o}$（字母 $o$），疑为 $C^{0}$（数字零）之误，此处径改；下文 (5.10) 同。——译者注}
\]
其中常数 $C\leq128$。
\end{theorem}

这个陈述表达的是 $E$ 在 $C^{0}$-\trTopology{}下于\trIdentityMap{}处的连续性。%
\cref{thm:c0-energy-estimate} 的证明也将推迟到下文的 5.6 节进行。

\begin{proposition}\label{prop:energy-properties}
\trEnergy{}函数 $E:\mathcal{D}\to\R$ 满足下列诸性质：
\begin{align*}
\text{(i)}\qquad& E(\varphi)\geq0,\ \text{且}\ E(\varphi)=0%
\Longleftrightarrow\varphi=\id\\
\text{(ii)}\qquad& E(\varphi)\ =\ E\bigl(\varphi^{-1}\bigr)\\
\text{(iii)}\qquad& E\bigl(\vartheta\varphi\vartheta^{-1}\bigr)=E(\varphi)\\
\text{(iv)}\qquad& E(\varphi\psi)\leq E(\varphi)+E(\psi),
\end{align*}
其中 $\varphi,\psi\in\mathcal{D}$，而 $\vartheta$ 是 $\R^{2n}$ 的一个%
\trSymplecticDiffeomorphism。
\end{proposition}

\begin{proof}
利用后面将要证明的非平凡事实 \eqref{eq:energy-positive}，(i) 与 (ii) 的%
证明由 $E$ 的定义立得。显然
\[
l\bigl(\varphi_{H}^{[0,1]}\bigr)=l\bigl(\varphi_{\overline{H}}^{[0,1]}\bigr)
\]
蕴含 $E(\varphi)=E(\varphi^{-1})$。由
\[
\vartheta\circ\varphi_{H}^{t}\circ\vartheta^{-1}=\varphi_{H_{\vartheta}}^{t}
\]
我们得到
\[
l\bigl(\bigl(\vartheta\circ\varphi_{H}\circ\vartheta^{-1}\bigr)^{[0,1]}\bigr)%
=l\bigl(\varphi_{H_{\vartheta}}^{[0,1]}\bigr)
\]
蕴含 (iii)。最后，由\cref{prop:calculus-formulae}我们有%
$\varphi_{H}^{t}\circ\varphi_{K}^{t}=\varphi_{H\#K}^{t}$，且
\begin{align*}
l\bigl(\varphi_{H\#K}^{[0,1]}\bigr)\ &=\ \int_{0}^{1}\Bigl\{\sup_{x}%
(H\#K)_{t}-\inf_{x}(H\#K)_{t}\Bigr\}\,dt\\
&\leq\ \int_{0}^{1}\Bigl(\sup_{x}H_{t}-\inf_{x}H_{t}\Bigr)\,dt+\int_{0}^{1}%
\Bigl(\sup_{x}K_{t}-\inf_{x}K_{t}\Bigr)\,dt\\
&=\ l\bigl(\varphi_{H}^{[0,1]}\bigr)+l\bigl(\varphi_{K}^{[0,1]}\bigr)
\end{align*}
蕴含 $E(\varphi_{H}\circ\varphi_{K})\leq E(\varphi_{H})+E(\varphi_{K})$，%
最后的论断 (iv) 随之得证。
\end{proof}

我们能对 $\mathcal{D}$ 中的\textbf{\trCommutator}（commutator）导出一个\trEstimate；%
\trCommutator{}记作
\begin{equation}\label{eq:commutator-def}
[\varphi,\psi]:=\varphi\psi\varphi^{-1}\psi^{-1},
\end{equation}
并且我们首先断言：
\begin{equation}\label{eq:commutator-energy}
E\bigl([\varphi,\psi]\bigr)\leq2\min\bigl\{E(\varphi),E(\psi)\bigr\}
\end{equation}
对 $\varphi,\psi\in\mathcal{D}$ 成立。的确，利用\cref{prop:energy-properties}%
的陈述 (ii)--(iv)，我们可以得出\trEstimate
\[%
E(\varphi\psi\varphi^{-1}\psi^{-1})\leq E(\varphi\psi\varphi^{-1})+E\bigl(\psi^{-1}\bigr)\leq2E(\psi)%
\]
以及%
$E(\varphi\psi\varphi^{-1}\psi^{-1})\leq E(\varphi)+E\bigl(\psi\varphi^{-1}\psi^{-1}\bigr)=2E(\varphi)$，%
论断即得证。现在，对那些\trSupport{}含于一个固定\trOpenSet{}内的映射，%
我们为其\trCommutator{}的\trEnergy{}推导一条「先验」\trEstimate。

\begin{proposition}\label{prop:commutator-support}
若 $U\subset\R^{2n}$ 开且\trBounded，且 $\vartheta\in\mathcal{D}$ 满足%
$\vartheta(U)\cap U=\varnothing$，则
\[
E\bigl([\varphi,\psi]\bigr)\leq4E(\vartheta)
\]
对一切使 $\operatorname{supp}(\varphi)$ 与 $\operatorname{supp}(\psi)$ 含于%
$U$ 之中的 $\varphi,\psi\in\mathcal{D}$ 成立。
\end{proposition}

\begin{proof}
定义
\[
\gamma:=\varphi\,\vartheta^{-1}\,\varphi^{-1}\,\vartheta\in\mathcal{D}
\]
则 $\gamma\big|_{U}=\varphi\big|_{U}$，因为 $\vartheta(U)\cap U=\varnothing$%
且 $\operatorname{supp}(\varphi)\subset U$。从而在 $U$ 上%
$\varphi\psi\varphi^{-1}\psi^{-1}=\gamma\psi\gamma^{-1}\psi^{-1}$；又由于这%
两个映射在 $U$ 之外都是\trIdentityMap，它们在 $\R^{2n}$ 上相等。于是，由%
\eqref{eq:commutator-energy} 与\cref{prop:energy-properties}，
\[
E\bigl([\varphi,\psi]\bigr)=E\bigl([\gamma,\psi]\bigr)\leq2E(\gamma)%
\leq2E\bigl(\vartheta^{-1}\bigr)+2E(\vartheta)=4E(\vartheta)
\]
此即所要证的。
\end{proof}

\cref{prop:energy-properties} 使我们能借助\trEnergy{}定义一个特定的\trMetric{}%
$d:\mathcal{D}\times\mathcal{D}\to[0,\infty]$：
\begin{equation}\label{eq:metric-d}
d(\varphi,\psi):=E\bigl(\varphi^{-1}\psi\bigr).
\end{equation}
特别地，$\psi\in\mathcal{D}$ 的\trEnergy{} $E(\psi)=d(\id,\psi)$ 就是它到%
\trIdentityMap{}的\trDistance；而且我们注意到，由\cref{prop:calculus-formulae}的最后一个%
公式，
\begin{equation}\label{eq:metric-inf}
d(\varphi,\psi)=\inf\Bigl\{\,\lVert H-K\rVert\ \Bigm|\ H\ \text{生成}\ %
\psi\ \text{且}\ K\ \text{生成}\ \varphi\,\Bigr\}.
\end{equation}

\begin{theorem}\label{thm:bi-invariant-metric}
函数 $d$ 是 $\mathcal{D}$ 上的一个\trBiInvariantMetric，即它对属于%
$\mathcal{D}$ 的一切 $\varphi,\psi$ 与 $\vartheta$ 满足：
\begin{align*}
\text{(i)}\qquad& d(\varphi,\psi)\geq0\ \text{且}\ d(\varphi,\psi)=0%
\Longleftrightarrow\varphi=\psi\\
\text{(ii)}\qquad& d(\varphi,\psi)\leq d(\varphi,\vartheta)+d(\vartheta,\psi)%
\ \text{且}\ d(\varphi,\psi)\ =\ d(\psi,\varphi)\\
\text{(iii)}\qquad& d(\vartheta\varphi,\vartheta\psi)=d(\varphi,\psi)%
=d(\varphi\vartheta,\psi\vartheta).
\end{align*}
\end{theorem}

\cref{thm:bi-invariant-metric} 是\cref{prop:energy-properties}的一个直接%
推论。在下文中，我们可以把 $(\mathcal{D},d)$ 当作一个\textbf{\trTopologicalGroup}%
（topological group）来看待。注意，只要诸\trSupport{}一致\trBounded，我们这个%
$\mathcal{D}$ 上的内蕴\trMetric{}关于 $C^{0}$-\trMetric{}就是连续的。的确，由%
\cref{thm:energy-capacity-inequality}与\cref{thm:c0-energy-estimate}，
\[
\sup\bigl\{\,c_{0}(U)\ \Bigm|\ U\ \text{开且}\ \psi(U)\cap\varphi(U)=\varnothing\,\bigr\}%
\ \leq\ d(\varphi,\psi)
\]
\begin{equation}\label{eq:d-c0-estimate}
d(\varphi,\psi)\leq128\cdot\operatorname{diameter}\bigl(\operatorname{supp}%
\,\varphi^{-1}\psi\bigr)\cdot\sup_{x}\bigl\lvert\varphi(x)-\psi(x)\bigr\rvert
\end{equation}
对属于 $\mathcal{D}$ 的每一对 $\varphi$ 与 $\psi$ 成立。此外，群 $\mathcal{D}$%
还容有一个\textbf{\trOrderStructure}（order structure）。定义 $\mathcal{D}$ 的正子集%
$\mathcal{D}^{+}\subset\mathcal{D}$ 为
\[
\mathcal{D}^{+}=\bigl\{\,\psi_{H}\ \bigm|\ H\in\mathcal{H}\ \text{且}\ %
H\geq0\,\bigr\}
\]
而 $\mathcal{D}$ 的负子集 $\mathcal{D}^{-}\subset\mathcal{D}$ 为
\[
\mathcal{D}^{-}=\bigl\{\,\psi_{H}\ \bigm|\ H\in\mathcal{H}\ \text{且}\ %
H\leq0\,\bigr\}.
\]
若 $\varphi\in\mathcal{D}^{+}$，则 $\varphi^{-1}\in\mathcal{D}^{-}$，并且对%
$\R^{2n}$ 的每个\trSymplecticDiffeomorphism{} $\vartheta$ 都有%
$\vartheta\varphi\vartheta^{-1}\in\mathcal{D}^{+}$。这由%
\cref{prop:calculus-formulae}可得。我们将在下文 5.4 节中证明
\begin{equation}\label{eq:positive-negative-inter}
\mathcal{D}^{+}\cap\mathcal{D}^{-}=\{\id\}.
\end{equation}
因此，正如我们将要证明的，可以如下在 $\mathcal{D}$ 上定义一个\textbf{\trPartialOrder}%
（partial order）结构 $\geq$：
\begin{equation}\label{eq:partial-order}
\varphi\geq\psi\quad\Longleftrightarrow\quad\varphi\circ\psi^{-1}%
\in\mathcal{D}^{+}.
\end{equation}

\begin{proposition}\label{prop:order-properties}
设 $\varphi,\psi$ 与 $\vartheta\in\mathcal{D}$。
\begin{align*}
\text{(i)}\qquad& \text{若}\ \varphi\geq\psi\ \text{且}\ \psi\geq\varphi\ %
\text{则}\ \varphi=\psi\\
\text{(ii)}\qquad& \text{若}\ \varphi\geq\psi\ \text{且}\ \psi\geq\vartheta\ %
\text{则}\ \varphi\geq\vartheta.
\end{align*}
此外，若 $\varphi\geq\psi$，则 $\varphi\vartheta\geq\psi\vartheta$。
\end{proposition}

\begin{proof}
若 $\varphi\circ\psi^{-1}\in\mathcal{D}^{+}$ 且 $\psi\circ\varphi^{-1}%
\in\mathcal{D}^{+}$，则 $\varphi\circ\psi^{-1}\in\mathcal{D}^{+}\cap\mathcal{D}^{-}$，%
从而由 \eqref{eq:positive-negative-inter} 得 $\varphi=\psi$，(i) 得证。其次%
注意到：若 $\varphi$ 与 $\psi$ 属于 $\mathcal{D}^{+}$，则由%
\cref{prop:calculus-formulae}，$\varphi\circ\psi\in\mathcal{D}^{+}$ 也成立，%
这就蕴含 (ii)。
\end{proof}

我们以属于 Y.~\trEliashberg{}与 L.~\trPolterovich{} [76] 的两个观察来结束%
本节。他们考察一个任意的双不变\textbf{\trPseudoMetric}（pseudo-metric）%
$\rho:\mathcal{D}\times\mathcal{D}\to[0,\infty)$。按定义，这个 $\rho$ 满足%
\cref{thm:bi-invariant-metric}的性质 (i)--(iii)，只是当 $\varphi\neq\psi$ 时%
并不要求 $\rho(\varphi,\psi)\neq0$。回顾 $E(\psi)=d(\id,\psi)$，并仿照%
\cref{prop:commutator-support}的证明进行论证，便发现
\begin{equation}\label{eq:pseudo-commutator}
\rho\bigl(\id,[\varphi,\psi]\bigr)\leq4\rho(\id,\vartheta)
\end{equation}
对一切满足 $\vartheta(U)\cap U=\varnothing$ 与%
$\operatorname{supp}(\varphi)$、$\operatorname{supp}(\psi)\subset U$（此处%
$U\subset\R^{2n}$ 为某个\trOpenSet）的 $\varphi,\psi$ 与 $\vartheta\in\mathcal{D}$%
成立。循着 H.~\trHofer{} [116]，对每个开的\trBounded{}子集%
$U\subset\R^{2n}$，人们可以联系其隶属于 $\rho$ 的所谓%
\trDisplacementEnergy{} $e_{\rho}(U)$，即定义
\begin{equation}\label{eq:displacement-energy}
e_{\rho}(U):=\inf\bigl\{\,\rho(\id,\vartheta)\ \bigm|\ \vartheta\in\mathcal{D}\
\text{满足}\ \vartheta(U)\cap U=\varnothing\,\bigr\}.
\end{equation}
于是由 \eqref{eq:pseudo-commutator} 得
\[
4e_{\rho}(U)\ \geq\ \rho\bigl(\id,[\varphi,\psi]\bigr)
\]
对每一对其\trSupport{}含于 $U$ 内的 $\varphi,\psi\in\mathcal{D}$ 成立。若%
$\rho$ 不仅是一个\trPseudoMetric{}而且还是一个\trMetric，则只要 $\varphi$ 与 $\psi$ 不可%
交换，便有 $\rho\bigl(\id,[\varphi,\psi]\bigr)>0$。从而 $e_{\rho}(U)>0$，%
我们就证明了

\begin{proposition}\label{prop:displacement-positive}
若 $\rho$ 是 $\mathcal{D}$ 上的一个\trBiInvariantMetric，则%
\trDisplacementEnergy{} $e_{\rho}$ 在\trOpenSet{}上取正值。
\end{proposition}

这特别适用于我们那个特定的\trBiInvariantMetric{} $\rho=d$。然而要注意，%
这个 $d$ 之为\trMetric{}的特性是由\cref{thm:energy-capacity-inequality}中的%
非平凡\trEstimate{}来断定的，由它可知：对每个\trOpenSet{} $U$ 都有%
$e_{d}(U)\geq c_{0}(U)>0$。特殊的\trDisplacementEnergy{} $e_{d}$ 将在后文中%
更详细地研究。

为了构造 $\mathcal{D}$ 上的内蕴\trBiInvariantMetric{} $d$，我们从%
$C^{\infty}_{c}\bigl(\R^{2n}\bigr)$ 的上确界\trNorm{} $\lVert H\rVert$ 出发。%
当然，我们同样可以从 $L^{p}$-\trNorm
\[
\lVert H\rVert_{p}:=\Bigl(\int\ \bigl\lvert H(x)\bigr\rvert^{p}\ dx\Bigr)^{\frac{1}{p}}
\]
出发，它在 $\mathcal{D}$ 之下显然\trInvariant，即当 $\varphi\in\mathcal{D}$ 时%
$\lVert H\circ\varphi\rVert_{p}=\lVert H\rVert_{p}$。于是我们得到%
$d_{p}(\varphi,\psi)=E_{p}(\varphi\psi^{-1})$\footnote{原书印作「$E_{p}(\varphi\psi^{-1})$」，与 (5.8) 之 $E(\varphi^{-1}\psi)$ 逆序不一致，此处照录原书面貌。——译者注}，它定义了 $\mathcal{D}$ 上的%
一个双不变\trPseudoMetric。然而结果表明，这个 $d_{p}$ 当 $1\leq p<\infty$ 时并不是%
\trMetric。这一事实最近由 Y.~\trEliashberg{}与 L.~\trPolterovich{} [76] 证明。%
他们更精确地证明了：
\[
E_{p}(\psi)=0,\quad\text{若}\ p=1,
\]
它对每个\trCalabi{}不变量为零的 $\psi\in\mathcal{D}$ 成立，即
\[
\int_{0}^{1}\int_{\R^{2n}}H(t,x)\ dx\ dt=0
\]
当 $\psi=\psi_{H}$（$H\in\mathcal{H}$）时成立。此外，当 $1<p<\infty$ 时，对一切%
$\psi\in\mathcal{D}$ 都有 $E_{p}(\psi)=0$。
